\subsection{Positive elements}

Let E be a complex Hilbert space. Let u be an endomorphism of E. Recall (EVT, V, p.~45, def. 6) that u is said to be positive if one has \(\langle x|u(x)\rangle\geqslant0\) for all \(x\in E\) . The endomorphism u is then Hermitian (loc. cit.). Moreover, if F is a complex Hilbert space and if \(v\in\mathcal{L}(\mathrm{F};\mathrm{E})\) , then the endomorphism \(v^{*}uv\) of F is positive (EVT, V, p.~45, prop. 12).

PROPOSITION 8. --- Let E be a complex Hilbert space. Let u be an endomorphism of E. The following conditions are equivalent:

\begin{enumerate}
\def\labelenumi{(\roman{enumi})}
\item
  The endomorphism u is positive;
\item
  The image of u under the numerical map is contained in \(\mathbf{R}_{+}\) ;
\item
  The endomorphism \(u\) is a positive element of the stellar algebra \(\mathcal{L}(\mathrm{E})\) ;
\item
  There exists a Hermitian element v of \(\mathcal{L}(\mathrm{E})\) such that \(u = v^2\) ;
\item
  There exists a continuous linear map v from E into a complex Hilbert space F such that \(u = v^{*}v\) .
\end{enumerate}

By EVT, V, p.~45, def. 6, u is positive if and only if it is Hermitian and if \(\langle x|u(x)\rangle\geqslant0\) for all \(x\in E\) . (i) The implication (i) \(\Longrightarrow\) (ii) therefore follows from the definition of the numerical image.

\begin{enumerate}
\def\labelenumi{(\roman{enumi})}
\setcounter{enumi}{1}
\item
  \(\Longrightarrow\) (iii): the hypothesis implies that u is Hermitian (EVT, V, p.~45, and remark, p.~2) and its spectrum is contained in \(R_{+}\) (prop. 6); consequently, u is a positive element of the stellar algebra \(\mathcal{L}(\mathrm{E})\) .
\item
  \(\Longrightarrow\) (iv): this is a special case of prop. 16 of I, p.~118.
\item
  \(\Longrightarrow\) (v) is immediate.
\item
  \(\Longrightarrow\) (i): let F be a complex Hilbert space and let \(v \in \mathcal{L}(\mathrm{E};\mathrm{F})\) such that \(u = v^{*}v\) . Let \(x \in E\) . We have \(\langle x|u(x)\rangle = \langle x|(v^{*}v)(x)\rangle = \|v(x)\|^{2}\) , which proves that u is positive.
\end{enumerate}

Recall (EVT, V, p.~45, remark 1) that, for every hermitian element u of \(\mathcal{L}(\mathrm{E})\) , we set

\[
m(u) = \inf_{\substack{x\in \mathrm{E}\\ \| x\| = 1}}\langle x|u(x)\rangle = \inf \iota (u) = \inf_{x\in \operatorname{E}\text{-}\{0\}}\frac{\langle x|u(x)\rangle}{\|x\|^ {2}},
\]

\[
\operatorname{M}(u) = \sup_{\substack{x\in \operatorname{E}\\ \| x\| = 1}}\langle x|u(x)\rangle = \sup \iota (u) = \sup_{x\in \operatorname{E}\text{-}\{0\}}\frac{\langle x|u(x)\rangle}{\|x\|^ {2}}.
\]

\[
\text{If } \mathrm{E} = \{0\}\text{, then } \mathrm{M}(u) = -\infty,\ m(u) = +\infty\text{, and } \iota(u) = \varnothing.
\]

Suppose E is nonzero; we then have \(m(u) \leqslant \mathrm{M}(u)\) and the numerical range of u is an interval with endpoints \(m(u)\) and \(\mathrm{M}(u)\) . By prop. 6, \(\mathrm{Sp}(u)\) is contained in the interval \([m(u), \mathrm{M}(u)]\) . More precisely:

PROPOSITION 9. --- Let E be a complex Hilbert space and let u be a hermitian element of \(\mathcal{L}(\mathrm{E})\) .

\begin{enumerate}
\def\labelenumi{\alph{enumi})}
\item
  We have \(m(u) = \inf \operatorname{Sp}(u)\) and \(\operatorname{M}(u) = \sup \operatorname{Sp}(u)\) ;
\item
  If E is not zero, we have \(\|u\|=\sup(|m(u)|,|\mathrm{M}(u)|)\) .
\end{enumerate}

Let \(\lambda \in \mathbf{R}\) . For \(\lambda\) to be a lower bound of the spectrum of \(u\) , it is necessary and sufficient that \(u - \lambda \geqslant 0\) . This is equivalent (prop. 8, (ii)) to the condition \(\langle x|u(x)\rangle \geqslant \lambda \|x\|^2\) for all \(x \in E\) , that is, to \(m(u) \geqslant \lambda\) . This proves that \(m(u)\) is the lower bound of \(\mathrm{Sp}(u)\) . Similarly, one verifies that \(\mathrm{M}(u)\) is the upper bound of \(\mathrm{Sp}(u)\) .

Since u is normal, we have \(\varrho(u)=\|u\|\) (cor. 1 of I, p.~108). Since \(E\neq\{0\}\) , the spectrum of u is not empty (cor. 1 of I, p.~26) and \(\varrho(u)\) is the radius of the smallest disk with center 0 that contains \(\mathrm{Sp}(u)\) (th. 1 of I, p.~24), hence b) follows from a).

